\section{Langkah-Langkah Pemecahan Masalah}

Pengembangan Judol Detector diawali dari identifikasi masalah utama:
\textbf{bagaimana mendeteksi indikasi konten judi online pada halaman web
secara otomatis dan efisien?} Masalah ini dipecah menjadi lima langkah
sistematis berikut.

\begin{enumerate}
    \item \textbf{Identifikasi pola dan variasi teks.}
    Konten judi online tidak selalu ditulis secara eksplisit. Selain kata
    kunci tetap seperti \texttt{slot}, \texttt{gacor}, dan \texttt{jackpot},
    terdapat pola khusus (misalnya \texttt{SLOT88}, \texttt{MAXWIN234}) serta
    variasi penulisan dengan karakter mirip secara visual (misalnya
    \texttt{G4C0R}, \texttt{H0KI88}).

    \item \textbf{Pemilihan strategi pencocokan string.}
    Untuk kata kunci tetap, dipilih algoritma exact matching: KMP dan
    Boyer-Moore, Rabin-Karp, dan Aho-Corasick. Untuk pola dengan variasi visual, digunakan Regular Expression
    dan Weighted Levenshtein Distance.

    \item \textbf{Perancangan mekanisme pemindaian halaman.}
    Sistem harus mampu membaca teks dari DOM halaman web, mengabaikan elemen
    yang tidak relevan (\texttt{script}, \texttt{style}, \texttt{input}), serta
    memproses teks dari gambar melalui OCR.

    \item \textbf{Penyajian hasil kepada pengguna.}
    Hasil deteksi perlu ditampilkan secara intuitif melalui highlight pada
    elemen DOM, tooltip saat hover, dan ringkasan statistik pada popup
    extension.

    \item \textbf{Optimasi performa.}
    Untuk menghindari pemrosesan berulang, diterapkan mekanisme deduplikasi
    hasil antar-algoritma dan penyimpanan cache pada pemindaian teks dan OCR.
\end{enumerate}

\section{Pemetaan Masalah menjadi Elemen Algoritma}

Setiap aspek masalah yang telah diidentifikasi dipetakan ke elemen algoritma
atau modul yang bertanggung jawab menyelesaikannya. Pemetaan ini
memperjelas peran masing-masing algoritma dalam sistem secara keseluruhan.

\begin{table}[H]
\centering
\caption{Pemetaan masalah ke elemen algoritma}
\label{tab:pemetaan-algoritma}
\begin{tabular}{|p{4.5cm}|p{8.5cm}|}
\hline
\textbf{Aspek Masalah} & \textbf{Elemen Algoritma / Modul} \\
\hline
Kata kunci tetap dari \texttt{keyword.txt} &
KMP (failure function), Boyer-Moore (bad character table), Rabin-Karp
(rolling hash), Aho-Corasick (trie \& failure links) \\
\hline
Pola \texttt{<kata><angka>} (contoh: \texttt{MAXWIN234}) &
RegEx engine JavaScript dengan pola \texttt{[\textbackslash{}p\{L\}]\{2,\}[0-9]\{2,3\}} \\
\hline
Variasi penulisan visual (contoh: \texttt{G4C0R}, \texttt{H0KI88}) &
Weighted Levenshtein Distance dengan biaya substitusi karakter mirip
yang lebih rendah \\
\hline
Pemrosesan teks dari gambar & OCR Tesseract.js \(\rightarrow\) pipeline
pencocokan string \\
\hline
Penggabungan hasil berbagai algoritma & Modul \texttt{scanner.ts}:
normalisasi, deduplikasi, cache LRU \\
\hline
Penandaan hasil pada halaman web & DOM manipulation: wrapper
\texttt{<span>}, tooltip, blur \\
\hline
\end{tabular}
\end{table}

\textbf{Pemetaan ke KMP.}
Masalah yang dihadapi adalah mencari kemunculan beberapa pola (keyword) dalam
teks yang panjang dan dinamis (DOM halaman web). KMP dipilih karena setiap
kali terjadi ketidakcocokan, informasi yang sudah cocok sebelumnya tidak
dibuang. Failure function memungkinkan pola digeser tepat ke posisi di mana
prefiks yang sudah cocok bisa dimanfaatkan kembali. Hal ini sangat cocok
untuk teks halaman web yang seringkali mengandung pengulangan huruf atau
frasa, sehingga jumlah perbandingan karakter tetap linear.

\textbf{Pemetaan ke Boyer-Moore.}
BM dipilih karena keunggulannya pada kasus teks alfanumerik yang umum di
halaman web. Ketika karakter pada teks tidak ditemukan dalam pola sama
sekali, BM dapat melompati sebagian besar teks tanpa perlu membandingkan
karakter satu per satu. Karakteristik ini sangat menguntungkan pada
pencarian keyword pendek (misalnya \texttt{slot}, \texttt{rtp}) dalam teks
halaman yang panjang, karena jumlah pergeseran bisa sangat besar.

\textbf{Pemetaan ke Rabin-Karp.}
RK dipilih karena kemampuannya mencari banyak pola sekaligus melalui
pendekatan rolling hash. Setiap keyword dihitung nilai hash-nya terlebih
dahulu; kemudian window hash pada teks digeser satu karakter per langkah.
Ketika nilai hash window cocok dengan hash pola, verifikasi karakter
dilakukan untuk menghindari hash collision. Pendekatan ini efisien ketika
banyak keyword pendek perlu ditemukan dalam teks yang sama karena
perbandingan hash jauh lebih cepat daripada perbandingan karakter penuh.

\textbf{Pemetaan ke Aho-Corasick.}
Aho-Corasick dipilih untuk skenario pencarian banyak keyword secara
simultan. Semua keyword dikompilasi menjadi trie dengan failure links,
sehingga teks hanya perlu dilintasi satu kali untuk menemukan seluruh
kemunculan semua keyword sekaligus. Berbeda dengan KMP dan BM yang
memproses satu keyword per lintasan, Aho-Corasick menyelesaikan seluruh
daftar keyword dalam satu traversal linear, menjadikannya sangat efisien
untuk keyword list yang panjang.

\section{Fitur Fungsional dan Arsitektur Extension}

\subsection{Fitur Fungsional}

Judol Detector menyediakan fitur-fungsional berikut kepada pengguna.

\begin{itemize}
    \item \textbf{Pemindaian teks halaman.}
    Content script membaca seluruh teks yang tampil pada DOM menggunakan
    \texttt{TreeWalker}, mengabaikan elemen \texttt{script},
    \texttt{style}, \texttt{noscript}, \texttt{textarea}, dan
    \texttt{input} agar pemindaian fokus pada konten yang benar-benar
    dibaca pengguna.

    \item \textbf{Pemindaian multi-algoritma.}
    Enam algoritma dijalankan secara independen pada setiap node teks:
    KMP, Boyer-Moore, Rabin-Karp, Aho-Corasick, Regular Expression, dan
    Weighted Levenshtein Distance.

    \item \textbf{Highlight dan tooltip.}
    Teks yang terdeteksi dibungkus dalam elemen \texttt{<span>} dengan
    class khusus. Tooltip muncul saat pengguna mengarahkan kursor,
    menampilkan keyword, algoritma asal, jumlah kemunculan, dan waktu
    eksekusi.

    \item \textbf{Blur otomatis.}
    Pengguna dapat mengaktifkan efek blur pada teks terdeteksi melalui
    popup extension. Gambar hasil OCR yang dinyatakan positif akan
    diburamkan secara otomatis.

    \item \textbf{OCR gambar.}
    Gambar yang terlihat pada viewport diproses langsung menggunakan
    Tesseract.js. Sistem tidak menggunakan metadata gambar seperti nama file,
    \texttt{alt}, \texttt{aria-label}, atau teks container sebagai sumber
    deteksi, sehingga keputusan positif berasal dari teks yang berhasil
    diekstraksi melalui OCR.

    \item \textbf{Rescan dan cache.}
    Tombol rescan pada popup menghapus highlight lama dan menjalankan
    pipeline deteksi dari awal. Hasil pemindaian disimpan dalam cache LRU
    agar pemrosesan berulang pada teks yang sama menjadi instan.
\end{itemize}

\subsection{Arsitektur Extension}

Sistem dibangun dengan tiga komponen utama yang saling berkomunikasi.

\begin{enumerate}
    \item \textbf{Content Script (\texttt{src/content/}).}
    Bertugas menelusuri DOM, menjalankan \texttt{scanTextForJudol},
    memasang highlight, tooltip, efek blur, serta menangani pesan dari
    popup melalui Chrome Messaging API.

    \item \textbf{Popup (\texttt{src/popup/}).}
    Antarmuka berbasis React yang menampilkan statistik real-time:
    total match, waktu eksekusi per algoritma, distribusi keyword, dan
    kontrol toggle (blur, OCR).

    \item \textbf{Modul Algoritma (\texttt{src/algorithms/}).}
    Kumpulan fungsi murni yang menerima \texttt{(text, keywords)} dan
    mengembalikan \texttt{AlgorithmResult}. Modul ini tidak bergantung
    pada API browser sehingga dapat diuji secara mandiri melalui unit
    test.
\end{enumerate}

\section{Contoh Ilustrasi Kasus}

Berikut adalah ilustrasi bagaimana pipeline deteksi bekerja pada sebuah
potongan teks halaman web.

\textbf{Teks halaman:}
\begin{verbatim}
Daftar sekarang di GACOR99 dan dapatkan MAXWIN234!
Main slot gacor hari ini. RTP tinggi, jackpot menanti.
\end{verbatim}

\textbf{Daftar keyword (\texttt{keyword.txt}):} \texttt{slot}, \texttt{gacor},
\texttt{jackpot}, \texttt{rtp}, \texttt{bonus deposit}

\textbf{Langkah deteksi:}
\begin{enumerate}
    \item \textbf{KMP} menemukan \texttt{"slot"} pada indeks 56--59,
    \texttt{"gacor"} pada indeks 61--65, \texttt{"rtp"} pada indeks 77--79,
    dan \texttt{"jackpot"} pada indeks 89--95.
    \item \textbf{Boyer-Moore} menemukan hasil yang sama pada posisi
    identik untuk keempat keyword tersebut.
    \item \textbf{Rabin-Karp} menghasilkan hash match untuk \texttt{"slot"},
    \texttt{"gacor"}, \texttt{"rtp"}, dan \texttt{"jackpot"}, lalu
    memverifikasi karakter demi karakter hingga dinyatakan cocok.
    \item \textbf{Aho-Corasick} menelusuri teks satu kali menggunakan trie
    yang telah dikompilasi dari seluruh keyword, dan menemukan keempat
    keyword yang sama pada posisi yang identik.
    \item \textbf{RegEx} menemukan pola \texttt{GACOR99} (indeks 19--25)
    dan \texttt{MAXWIN234} (indeks 40--48) karena keduanya terdiri atas
    minimal dua huruf dan diikuti dua hingga tiga angka tanpa spasi.
    \item \textbf{Weighted Levenshtein} tidak menemukan kecocokan baru
    karena semua keyword pada daftar sudah terdeteksi secara eksak oleh
    algoritma sebelumnya.
\end{enumerate}

\textbf{Hasil setelah deduplikasi:}
\begin{itemize}
    \item \texttt{GACOR99} --- deteksi oleh RegEx
    \item \texttt{MAXWIN234} --- deteksi oleh RegEx
    \item \texttt{slot} --- deteksi oleh KMP / Boyer-Moore / Rabin-Karp / Aho-Corasick
    (diambil satu karena rentang sama)
    \item \texttt{gacor} --- deteksi oleh KMP / Boyer-Moore / Rabin-Karp / Aho-Corasick
    (diambil satu karena rentang sama)
    \item \texttt{rtp} --- deteksi oleh KMP / Boyer-Moore / Rabin-Karp / Aho-Corasick
    (diambil satu karena rentang sama)
    \item \texttt{jackpot} --- deteksi oleh KMP / Boyer-Moore / Rabin-Karp / Aho-Corasick
    (diambil satu karena rentang sama)
\end{itemize}

Setelah deduplikasi, keenam hasil tersebut di-highlight pada DOM dan
ringkasan statistiknya ditampilkan pada popup extension.
